%!TEX root = main.tex
\section{Marco Teórico}

\subsection{Filtrado Espacial y Suavizado de Imágenes}
El suavizado de imágenes es una técnica de filtrado espacial utilizada para reducir el ruido y los detalles irrelevantes de una imagen. El método más común es el filtro de caja (box filter) de $3 \times 3$, el cual asigna a cada píxel de salida el promedio de los valores de sus vecinos inmediatos en una ventana cuadrada. Matemáticamente, para un píxel en la posición $(x,y)$, el valor suavizado $S(x,y)$ se calcula como:

\begin{equation}
S(x,y) = \frac{1}{9} \sum_{i=-1}^{1} \sum_{j=-1}^{1} I(x+i, y+j)
\end{equation}

Donde $I$ es la imagen de entrada. Esta operación es una convolución discreta que, a diferencia de otros filtros más complejos, presenta una alta regularidad en los accesos a memoria, lo que la hace ideal para estudiar el impacto de la aceleración mediante hardware paralelo.

\subsection{Paralelismo con OpenMP}
OpenMP (Open Multi-Processing) es una interfaz de programación de aplicaciones (API) que admite la programación multiproceso de memoria compartida. En este proyecto, se exploran dos niveles de paralelismo ofrecidos por el estándar:

\begin{itemize}
    \item \textbf{Multicore CPU:} Utilizando la directiva \texttt{\#pragma omp parallel for}, el trabajo de iteración sobre las filas de la imagen se divide entre múltiples hilos de ejecución concurrentes, aprovechando todos los núcleos físicos del procesador.
    \item \textbf{Heterogeneous Offloading (OpenMP Target):} Permite delegar regiones de código a aceleradores externos. Es un enfoque de "alto nivel" que prioriza la portabilidad entre diferentes tipos de GPUs y FPGAs.
    \item \textbf{CUDA (Compute Unified Device Architecture):} Es una arquitectura de computación paralela de bajo nivel desarrollada por NVIDIA. A diferencia de OpenMP, CUDA ofrece un control granular sobre la jerarquía de memoria y la ejecución en la GPU, permitiendo optimizaciones extremas mediante el uso de mallas de bloques y hilos (\textit{Grids, Blocks and Threads}).
\end{itemize}

\subsection{Modelo de Ejecución en la GPU}
Tanto en OpenMP Target como en CUDA, la GPU ejecuta miles de hilos de forma concurrente. En CUDA, estos hilos se agrupan en bloques, y los bloques en una malla global. Esta arquitectura permite que el suavizado de un millón de píxeles se realice en paralelo, donde cada hilo se encarga de un solo dato, eliminando casi por completo la secuencialidad del algoritmo original.

\subsection{Análisis de Rendimiento}
Para la evaluación de las implementaciones, se utilizan métricas de tiempo de ejecución y análisis de "overhead". El overhead se define como el tiempo adicional requerido para la creación de hilos, sincronización y transferencia de datos a través del bus PCIe (en el caso de la GPU), lo cual puede superar el beneficio del paralelismo en conjuntos de datos pequeños.

\subsection{Formatos de Imagen PGM}
Se emplea el formato \textit{Portable GrayMap} (PGM) en su variante binaria P5. Este formato permite una representación eficiente de la matriz de intensidades de gris para resoluciones masivas de 8K ($8192 \times 8192$), facilitando la validación visual del suavizado mediante la inspección de los archivos de salida procesados por cada una de las versiones del kernel.